\documentclass[aspectratio=169,11pt]{beamer}

\usetheme{Madrid}
\usecolortheme{default}
\setbeamertemplate{navigation symbols}{}
\setbeamertemplate{footline}[frame number]

\usepackage{amsmath,amssymb,bm}
\usepackage{booktabs}
\usepackage{physics}
\usepackage{hyperref}
\usepackage{CJKutf8}
\usepackage[backend=biber,style=numeric,sorting=none]{biblatex}
\addbibresource{main.bib}

\title[\texorpdfstring{$^{171}$Yb Förster resonance}{171Yb Forster resonance}]%
{Förster resonance 对 \texorpdfstring{$^{171}\mathrm{Yb}$}{171Yb} UV CZ time-optimal 序列 fidelity 的影响}
\subtitle{Notebook-first 研究计划：no-jump 优化序列的完整 Lindblad 评估}
\author{Noise-Tolerant Quantum Control Optimization}
\date{\today}

\newcommand{\ketb}{\ket{b}}
\newcommand{\ketloss}{\ket{\mathrm{loss}}}
\newcommand{\Dforster}{\Delta_F}
\newcommand{\Fproj}{F_{\mathrm{proj}}}

\begin{document}
\begin{CJK*}{UTF8}{gbsn}

\begin{frame}
  \titlepage
\end{frame}

\begin{frame}{本阶段目标}
  \begin{block}{核心问题}
    既有 UV edge scan 使用 Rydberg decay 的 no-jump 非厄密有效生成元训练
    phase-only time-optimal 序列。现在的问题是：若加入 Peper et al. 指出的
    nearby Förster pair state，这些序列的实际 Lindblad fidelity 会下降多少？
  \end{block}

  \begin{itemize}
    \item 训练阶段：暂不重新优化，沿用已有 artifact 中的相位序列。
    \item 评估阶段：显式加入 Förster pair state，并用完整 Lindblad 主方程传播。
    \item 输出目标：扫描 \(\Dforster\) 与 coupling \(W\)，判断 resonance dip 和 edge-time sensitivity。
  \end{itemize}
\end{frame}

\begin{frame}{文献动机}
  \begin{itemize}
    \item Ma et al. 2022 首次在 \({}^{171}\mathrm{Yb}\) nuclear-spin qubit 上展示 Rydberg blockade gate，并指出强 blockade 可能和 hyperfine-induced Förster resonance 有关 \cite{maUniversalGateOperations2022}.
    \item Ma et al. 2023 使用 metastable qubit 与 302 nm 单光子 Rydberg excitation，实现 CZ fidelity \(0.980(1)\)，但误差预算仍有未解释项 \cite{maHighfidelityGatesMidcircuit2023a}.
    \item Peper et al. 2025 用 MQDT 找到此前 \(F=3/2\) Rydberg state 的 anomalous Förster resonance，并通过选择更干净的 \(F=1/2\) state 把 CZ fidelity 提升到 \(0.994(1)\) \cite{peperSpectroscopyModelingYb2025}.
  \end{itemize}
\end{frame}

\begin{frame}{为什么不能只用 no-jump 生成元}
  no-jump 近似对应只传播未发生 quantum jump 的 Kraus branch：
  \[
    G=-iH-\frac12\sum_j C_j^\dagger C_j .
  \]

  对只关心 Rydberg lifetime loss 的优化，这可以作为高效训练目标。
  但 Förster resonance 引入的是 coherent pair-state mixing：
  \[
    \ket{rr}\leftrightarrow \ketb .
  \]

  若 \(\ketb\) 有自己的 decay/leakage 路径，jump branch 也属于实际 trace-preserving channel 的一部分。
  因此评估 fidelity 时应使用完整 Lindblad 方程：
  \[
    \dot\rho=-i[H,\rho]+\sum_j C_j\rho C_j^\dagger
    -\frac12\sum_j\{C_j^\dagger C_j,\rho\}.
  \]
\end{frame}

\begin{frame}{第一阶段最小模型}
  Hilbert basis 在原 reduced UV basis 上增加一个 pair state 与一个 absorbing loss state：
  \[
  \{\ket{00},\ket{0c},\ket{0r},\ket{cc},\ket{W_{cr}},\ket{rr},\ketb,\ketloss\}.
  \]

  Hamiltonian：
  \[
    H(t)=H_0 + A(t)\left[\cos\phi(t)H_x+\sin\phi(t)H_y\right],
  \]
  \[
    H_0 = V\ket{rr}\bra{rr}
    +(V+\Dforster)\ketb\bra{b}
    +W(\ket{rr}\bra{b}+\ketb\bra{rr}).
  \]

  这里 \(\Dforster\) 是相对 target \(\ket{rr}\) branch 的 Förster defect。
\end{frame}

\begin{frame}{Decay channels}
  Rydberg decay 用 one-way Lindblad jumps 进入 \(\ketloss\)：
  \[
    C_{0r}=\sqrt{\gamma_r}\ketloss\bra{0r},
    \qquad
    C_{W}=\sqrt{\gamma_r}\ketloss\bra{W_{cr}},
  \]
  \[
    C_{rr}=\sqrt{2\gamma_r}\ketloss\bra{rr},
    \qquad
    C_b=\sqrt{2\gamma_r}\ketloss\bra{b}.
  \]

  第一阶段不加入 Doppler、laser phase noise、amplitude noise 或额外 \(m_F\) leakage。
  这些项可以在 resonance sensitivity 明确后逐项加入。
\end{frame}

\begin{frame}{Fidelity 口径}
  计算子空间代表态：
  \[
    \{\ket{00},\ket{0c},\ket{cc}\},
  \]
  权重 \((1,2,1)\) 对应完整 computational basis 中
  \(\ket{00},\ket{01}/\ket{10},\ket{11}\)。

  第一阶段使用 projected probe fidelity：
  \[
    \Fproj = \frac{1}{N}\sum_k
    \bra{\psi_k^{\mathrm{target}}}
    P_{\mathrm{active}}\rho_k(T)P_{\mathrm{active}}
    \ket{\psi_k^{\mathrm{target}}},
  \]
  其中 target phase 中的 local phases 数值优化。

  该口径比 no-jump diagonal process fidelity 更保守，因为 loss branch 不会被丢弃。
\end{frame}

\begin{frame}{Notebook 和执行路径}
  \begin{itemize}
    \item Notebook：
      \[
        \texttt{notebooks/yb171\_forster\_resonance.ipynb}
      \]
    \item 输入 artifact：
      \[
        \texttt{artifacts/v5/closed\_cr\_edge\_time\_optimal\_scan/}
      \]
      \[
        \texttt{rydberg\_decay\_65us\_dense\_selected\_phase\_rows.json}
      \]
    \item 默认先评估 \(\Omega_{\max}/2\pi=10\,\mathrm{MHz}\), edge \(=40\,\mathrm{ns}\) 的已优化序列。
    \item 扫描变量：\(\Dforster\in[-20,20]\,\mathrm{MHz}\)，\(W\in[0,5]\,\mathrm{MHz}\) 起步。
  \end{itemize}
\end{frame}

\begin{frame}{待验证结论}
  \begin{enumerate}
    \item \(W=0\) 时，完整 Lindblad projected fidelity 应接近已有 no-jump 序列的趋势，但数值口径不必完全相同。
    \item 当 \(|\Dforster|\) 为几 MHz 且 \(W\) 可观时，应检查 fidelity 是否出现 resonance dip。
    \item 比较不同 edge time 的同类序列，判断 Gaussian edge smoothing 是否降低 Förster sensitivity。
    \item 若 \(|\Dforster|\gg W\)，结果应接近 effective blockade-shift picture。
    \item 若 resonance dip 明显，后续优化应把 \(\ketb\) 显式纳入 GRAPE，而不是只使用非厄密有效生成元。
  \end{enumerate}
\end{frame}

\begin{frame}{参考文献}
  \printbibliography
\end{frame}

\end{CJK*}
\end{document}
